% Figure from MATH 451 notes, section 33. Compile with the notes preamble.
\begin{tikzpicture}
\begin{axis}[nfig, width=0.62\textwidth, height=0.38\textwidth,
  xmin=0, xmax=3.45, ymin=0, ymax=1.3,
  xtick={0.2,1.1,1.6,2.1,3.2}, xticklabels={$a$,$x_0 - \varepsilon$,$x_0$,$x_0 + \varepsilon$,$b$},
  ytick={0.525,1.05}, yticklabels={$\tfrac12 f(x_0)$,$f(x_0)$}, xlabel={}]
  \addplot[draw=none, fill=figB!12, domain=0.2:3.2, samples=100] {0.15+0.9*exp(-3*(x-1.6)^2)} \closedcycle;
  \addplot[draw=figR, thin, fill=figR!15] coordinates {(1.1,0) (2.1,0) (2.1,0.525) (1.1,0.525)} -- cycle;
  \addplot[black!40, dashed, thin] coordinates {(0,0.525) (1.1,0.525)};
  \addplot[black!40, dashed, thin] coordinates {(0,1.05) (1.6,1.05)};
  \addplot[figB, very thick, domain=0.2:3.2, samples=100] {0.15+0.9*exp(-3*(x-1.6)^2)};
  \addplot[only marks, mark=*, mark size=1.6pt, figB] coordinates {(1.6,1.05)};
  \node[font=\scriptsize, figR, align=center] at (axis cs:1.6,0.26) {area\\ $f(x_0)\,\varepsilon$};
  \node[font=\small, figB] at (axis cs:2.55,0.62) {$f$};
\end{axis}
\end{tikzpicture}
